\documentclass[11pt]{article}
\usepackage[margin=1in]{geometry}
\usepackage{amsmath,amssymb}
\usepackage{enumitem}
\title{MATH 1314 Precalculus Algebra: Section Summary}
\author{Timothy Trujillo}
\date{}
\begin{document}
\maketitle

\section*{Course Arc}
The course moves from algebraic fluency to functional thinking. Students practice solving equations, interpreting solution sets, reading graphs, composing and transforming functions, and using functions as models.

\section{Introduction to Precalculus}
\textbf{Objectives.} Connect practice, feedback, representation, and course objectives. Treat mathematics as a set of habits developed through revision.

\textbf{Example.} A student should be able to explain why the same relation might be described verbally, numerically, graphically, and symbolically.

\section{Equations and Inequalities}
\textbf{Objectives.}
\begin{itemize}[noitemsep]
  \item Plot and interpret ordered pairs in the coordinate plane.
  \item Solve equations equivalent to \(ax+b=cx+d\).
  \item Use \(i^2=-1\) in complex arithmetic.
  \item Solve quadratic equations by factoring, completing the square, and the quadratic formula.
  \item Represent inequalities using interval notation.
\end{itemize}
\textbf{Examples.}
\[
d=\sqrt{(x_2-x_1)^2+(y_2-y_1)^2}, \qquad
x=\frac{-b\pm\sqrt{b^2-4ac}}{2a}.
\]

\section{Functions}
\textbf{Objectives.}
\begin{itemize}[noitemsep]
  \item Evaluate and interpret \(f(a)\), \(f(x+h)\), and \(\frac{f(b)-f(a)}{b-a}\).
  \item Determine domains and ranges from formulas and graphs.
  \item Compute compositions \((f\circ g)(x)=f(g(x))\).
  \item Describe transformations \(a f(b(x-h))+k\).
  \item Find and verify inverse functions.
\end{itemize}
\textbf{Example.} For \(f(x)=\sqrt{x-3}\), the domain is \([3,\infty)\) because the radicand must be nonnegative.

\section{Linear Functions}
\textbf{Objectives.} Interpret slope, intercepts, point-slope form, slope-intercept form, and linear models.
\[
f(x)=mx+b.
\]
\textbf{Example.} If \(f(x)=35x+120\), then \(35\) is the rate of change and \(120\) is the initial value.

\section{Polynomials and Rational Functions}
\textbf{Objectives.}
\begin{itemize}[noitemsep]
  \item Analyze quadratic functions using vertex, intercepts, and concavity.
  \item Predict polynomial end behavior from degree and leading coefficient.
  \item Use division, factor, and remainder theorems.
  \item Analyze rational functions using zeros, holes, and asymptotes.
  \item Solve radical equations and check for extraneous solutions.
\end{itemize}
\textbf{Example.} The functions \(\frac{x^2-1}{x-1}\) and \(\frac{1}{x-1}\) both exclude \(x=1\), but the first has a removable discontinuity while the second has a vertical asymptote.

\section{Exponential and Logarithmic Functions}
\textbf{Objectives.}
\begin{itemize}[noitemsep]
  \item Interpret \(f(x)=ab^x\) using initial value and growth factor.
  \item Use \(\log_b(x)=y\) if and only if \(b^y=x\).
  \item Apply logarithm properties to expand, condense, and solve equations.
  \item Model growth and decay with exponential functions.
\end{itemize}
\textbf{Example.}
\[
2e^{0.4t}=10 \quad \Longrightarrow \quad t=\frac{\ln(5)}{0.4}.
\]

\section{Assessment}
Exam preparation should begin from learning objectives, continue through representative problem solving, and end with reflection on errors. The sample midterm and final are paired with recorded solution walkthroughs in the public archive.

\end{document}
